\chapter{Condições iniciais e constantes}\label{ap:dados}

As tabelas abaixo reproduzem o arquivo \arquivo{efemeride\_2026-01-01.json}, idêntico nas versões 2D e 3D. Os vetores de estado vêm do JPL Horizons (efeméride DE441), no instante 2026-01-01 0h TDB (JD 2461041,5), relativos ao baricentro do Sistema Solar, no referencial da eclíptica e equinócio de J2000. Os parâmetros gravitacionais vêm do cabeçalho da efeméride DE440 \cite{park2021}. Os planetas são os baricentros dos seus sistemas, e a ``Terra'' é o baricentro Terra--Lua.

\begin{table}[htb]
  \centering
  \caption{Parâmetros gravitacionais e raios}
  \label{tab:gm}
  \small
  \begin{tabular}{lcr}
    \toprule
    Corpo & $GM$ ($\mathrm{UA^3/dia^2}$) & Raio médio (km) \\
    \midrule
    Sol & $2{,}959122083\times10^{-4}$ & 695700,0 \\
Mercúrio & $4{,}912500195\times10^{-11}$ & 2439,4 \\
Vênus & $7{,}243452333\times10^{-10}$ & 6051,8 \\
Terra & $8{,}997011393\times10^{-10}$ & 6371,0 \\
Marte & $9{,}549548830\times10^{-11}$ & 3389,9 \\
Júpiter & $2{,}825345825\times10^{-7}$ & 69911,0 \\
Saturno & $8{,}459705993\times10^{-8}$ & 58232,0 \\
Urano & $1{,}292026565\times10^{-8}$ & 25362,0 \\
Netuno & $1{,}524357348\times10^{-8}$ & 24624,0 \\
    \bottomrule
  \end{tabular}
  \legend{O raio é o do corpo central (ex.: o da Terra, não o do baricentro Terra--Lua) e só é usado para o desenho. Fonte: JPL, via \arquivo{ferramentas/gerar\_efemeride.py}.}
\end{table}

\begin{table}[htb]
  \centering
  \caption{Posições iniciais (UA)}
  \label{tab:pos}
  \small
  \begin{tabular}{lrrr}
    \toprule
    Corpo & $x$ & $y$ & $z$ \\
    \midrule
    Sol & $-0{,}0030673162$ & $-0{,}0055333291$ & $+0{,}0001316842$ \\
Mercúrio & $-0{,}2182686205$ & $-0{,}4147409448$ & $-0{,}0135715854$ \\
Vênus & $+0{,}0858099334$ & $-0{,}7272957114$ & $-0{,}0149123724$ \\
Terra & $-0{,}1773370747$ & $+0{,}9622523453$ & $+0{,}0000748181$ \\
Marte & $+0{,}3375123607$ & $-1{,}3925353450$ & $-0{,}0372855426$ \\
Júpiter & $-1{,}6970703464$ & $+4{,}9233493856$ & $+0{,}0175579241$ \\
Saturno & $+9{,}5042760491$ & $+0{,}2522051287$ & $-0{,}3828038443$ \\
Urano & $+9{,}8772506590$ & $+16{,}7944820807$ & $-0{,}0655881149$ \\
Netuno & $+29{,}8690527247$ & $+0{,}5134061342$ & $-0{,}6989354910$ \\
    \bottomrule
  \end{tabular}
  \legend{Fonte: JPL Horizons.}
\end{table}

\begin{table}[htb]
  \centering
  \caption{Velocidades iniciais (UA/dia)}
  \label{tab:vel}
  \small
  \begin{tabular}{lrrr}
    \toprule
    Corpo & $v_x$ & $v_y$ & $v_z$ \\
    \midrule
    Sol & $7{,}17607143\times10^{-6}$ & $1{,}77813412\times10^{-7}$ & $-1{,}35885762\times10^{-7}$ \\
Mercúrio & $1{,}92390284\times10^{-2}$ & $-1{,}17374485\times10^{-2}$ & $-2{,}72329181\times10^{-3}$ \\
Vênus & $1{,}99457363\times10^{-2}$ & $2{,}40054371\times10^{-3}$ & $-1{,}11760768\times10^{-3}$ \\
Terra & $-1{,}72044191\times10^{-2}$ & $-3{,}11347577\times10^{-3}$ & $1{,}40487014\times10^{-7}$ \\
Marte & $1{,}41265880\times10^{-2}$ & $4{,}54038100\times10^{-3}$ & $-2{,}51221253\times10^{-4}$ \\
Júpiter & $-7{,}22343983\times10^{-3}$ & $-2{,}10136205\times10^{-3}$ & $1{,}70365041\times10^{-4}$ \\
Saturno & $-4{,}55908478\times10^{-4}$ & $5{,}56484997\times10^{-3}$ & $-7{,}86251448\times10^{-5}$ \\
Urano & $-3{,}41922498\times10^{-3}$ & $1{,}81055543\times10^{-3}$ & $5{,}10181776\times10^{-5}$ \\
Netuno & $-7{,}46211577\times10^{-5}$ & $3{,}15728516\times10^{-3}$ & $-6{,}32992143\times10^{-5}$ \\
    \bottomrule
  \end{tabular}
  \legend{Fonte: JPL Horizons.}
\end{table}

As demais constantes usadas são a unidade astronômica, $1\UA = 149\,597\,870{,}7\km$ (definição da IAU, 2012), e a obliquidade da eclíptica em J2000, $\varepsilon_0 = 23{,}4392911^\circ$.
